\subsection{Support of a Macaulay module}

PROPOSITION 2. — Let A be a Noetherian ring and M a finitely generated Macaulay A-module.

\begin{enumerate}
\def\labelenumi{\alph{enumi})}
\item
  The A-module M has no embedded associated prime ideals. \(^{1}\)
\item
  Let X be a closed irreducible subset of Supp(M) and Y a closed subset of X. We have
\end{enumerate}

\[
\operatorname{codim} (Y, X) + \operatorname{codim} (X, \operatorname{Supp} (M)) = \operatorname{codim} (Y, \operatorname{Supp} (M)).
\]

\begin{enumerate}
\def\labelenumi{\alph{enumi})}
\setcounter{enumi}{2}
\item
  The topological space Supp(M) is catenary (VIII, § 1, no. 2, def. 4).
\item
  Let \(X_{1}\) and \(X_{2}\) be irreducible components of Supp(M) and Y a closed subset of \(X_{1} \cap X_{2}\) . We have codim(Y, \(X_{1}\) ) = codim(Y, \(X_{2}\) ).
\item
  Let \(\mathfrak{p} \in \operatorname{Ass}(M)\) . One has \(\operatorname{prof}_{\mathrm{A}}(\mathfrak{p}; M) = 0\) ( \(\S\) 1, \(n^{\circ}\) 1, remark 2), hence \(\dim(M_{\mathfrak{p}}) = 0\) ( \(n^{\circ}\) 1, cor. of prop. 1), which implies that \(\mathfrak{p}\) is a minimal element of \(\operatorname{Supp}(M)\) .
\item
  Suppose first that Y is irreducible. Let p and q be the prime ideals of Supp(M) such that \(Y = V(\mathfrak{q})\) and \(X = V(\mathfrak{p})\) . It follows from prop. 1 that we have
\end{enumerate}

\[
\begin{array}{r l} \operatorname{codim} (Y, X) & = \dim (A _ {q} / p A _ {q}) = \dim (M _ {q}) - \dim (M _ {p}) \\ & = \operatorname{codim} (Y, \operatorname{Supp} (M)) - \operatorname{codim} (X, \operatorname{Supp} (M)). \end{array}
\]

The general case follows from VIII, § 1, no. 2, remark 3.

\begin{enumerate}
\def\labelenumi{\alph{enumi})}
\setcounter{enumi}{2}
\tightlist
\item
  Let X, Y, Z be closed irreducible subsets of Supp(M) such that \(Z \subset Y \subset X\) . The codimension of each of these subsets in Supp(M) is finite (VIII, § 1, no. 4, prop. 9 and § 3, no. 1, cor. 1 of prop. 2). We then deduce from b) the equality
\end{enumerate}

\[
\operatorname{codim} (Z, Y) + \operatorname{codim} (Y, X) = \operatorname{codim} (Z, X)
\]

which implies c) (VIII, § 1, no. 2, prop. 4).

\begin{enumerate}
\def\labelenumi{\alph{enumi})}
\setcounter{enumi}{3}
\tightlist
\item
  By b), we have \(\text{codim}(Y, X_1) = \text{codim}(Y, \text{Supp}(M)) = \text{codim}(Y, X_2)\) .
\end{enumerate}

In particular, if there exists a finitely generated Macaulay A-module M with support equal to Spec(A), the ring A is catenary and consequently every ring of fractions or every quotient ring of A is catenary (VIII, § 1, no. 3, remark 2).

Remark 1. — Under the hypotheses of prop. 2, it can happen that two irreducible components \(X_{1}\) and \(X_{2}\) of Supp(M) have an intersection Y reduced to a point and that we have \(\dim X_{1} \neq \dim X_{2}\) and \(\dim X_{2} \neq \text{codim}(Y, X_{2})\) (cf. exercise 4). However, this cannot happen when the ring A is local, as the following corollary shows.

COROLLARY. — Let A be a Noetherian local ring and M a nonzero finitely generated Macaulay A-module.

\begin{enumerate}
\def\labelenumi{\alph{enumi})}
\item
  All maximal chains of closed irreducible subsets of Supp(M) have length equal to dim(M).
\item
  For every closed subset X of Supp(M), we have
\end{enumerate}

\[
\operatorname{codim} (\mathrm{X}, \operatorname{Supp} (\mathrm{M})) = \dim (\operatorname{Supp} (\mathrm{M})) - \dim (\mathrm{X}).
\]

\begin{enumerate}
\def\labelenumi{\alph{enumi})}
\setcounter{enumi}{2}
\item
  All irreducible components of Supp(M) have the same dimension.
\item
  For every ideal J of A, we have
\end{enumerate}

\[
\operatorname{prof} _ {\mathrm{A}} (\mathrm{J}; \mathrm{M}) = \dim (\mathrm{M}) - \dim (\mathrm{M} / \mathrm{JM}).
\]

\begin{enumerate}
\def\labelenumi{\alph{enumi})}
\tightlist
\item
  A maximal chain of closed irreducible subsets of Supp(M) has as its smallest element \(\{\mathfrak{m}_{\mathrm{A}}\}\) and as its largest element an irreducible component X of
\end{enumerate}

Supp(M). Its length is equal to the codimension of \(\{m_{A}\}\) in X (prop. 2, c)); by prop. 2, b) applied to the closed subsets \(\{m_{A}\}\subset X\) , this is equal to \(\text{codim}(\{\mathfrak{m}_{\Lambda}\},\text{Supp}(M))\) , that is, to \(\dim(M)\) .

\begin{enumerate}
\def\labelenumi{\alph{enumi})}
\setcounter{enumi}{1}
\item
  This is a consequence of a) when the subset X is irreducible (VIII, § 1, no. 2, prop. 5); the general case follows from VIII, § 1, no. 1, prop. 1 and § 1, no. 2, remark 4.
\item
  This is a consequence of b).
\item
  We have \(\operatorname{prof}_{\mathrm{A}}(\mathrm{J};\mathrm{M}) = \operatorname{codim}(\operatorname{Supp}(\mathrm{M})\cap \mathrm{V}(\mathrm{J}),\operatorname{Supp}(\mathrm{M}))\) by the cor. of prop. 1 of no. 1. It suffices then to apply b) with \(\mathrm{X} = \operatorname{Supp}(\mathrm{M})\cap \mathrm{V}(\mathrm{J}) = \operatorname{Supp}(\mathrm{M} / \mathrm{JM})\) (II, § 4, no. 4, cor. of prop. 18).
\end{enumerate}

Remark 2.--- Let M be a finitely generated Macaulay A-module, and let p be an element of Supp(M). In view of th. 2 of § 1, no. 4, we have prof\_A(p; M) \textless{} +∞, and there exists an M-regular sequence of length prof\_A(p; M) consisting of elements of p. Let J denote the ideal of A generated by such a sequence; then the A-module M/JM is Macaulay (no. 1, example 3), and p is a minimal element of its support. Indeed, p contains J and hence belongs to the support of M/JM (II, § 4, no. 4, cor. of prop. 18); by corollary 1 of theorem 2 of § 1, no. 4, the ideal p is contained in an element of Ass(M/JM), but every prime ideal associated with a finitely generated Macaulay module is a minimal element of its support (prop. 2).
% semantic-fragments: legacy-input-seam
\relax{}
